\section{Related Work}
\label{sec:related_work}

\subsection{Edge-Native Network Intrusion Detection}
Deploying machine learning models directly on edge devices and industrial gateways has gained substantial traction due to strict latency and bandwidth budgets~\cite{sforzin2016rpids}. Recent works have investigated ensemble classifiers and meta-heuristic optimization on IoT single-board computers~\cite{baalaji2025ensemble,asghar2026hedid,zhang2026edgeics}. Similarly, deep autoencoders and generative architectures have been deployed on NVIDIA Jetson embedded platforms to detect industrial anomalies~\cite{hasan2025iiot,diazgorrin2026jetson}. However, existing research primarily optimizes offline model detection accuracy (such as F1-score or AUC-ROC) under idealized batch assumptions. In contrast, our work explicitly treats the \emph{runtime serving stack} as the primary object of scientific investigation, analyzing how precision representations and memory dispatch patterns impact operational throughput, cold-start latency, and execution parity.

\subsection{Serving Classical Machine Learning on Constrained Hardware}
While deep learning inference has benefited from extensive compiler optimizations, serving classical tree ensembles on embedded platforms remains challenging. The Open Neural Network Exchange (ONNX) specification standardizes operators such as \texttt{TreeEnsembleClassifier}, enabling models trained in frameworks like scikit-learn, XGBoost, or LightGBM to be executed across diverse backends. ONNX Runtime provides a highly optimized C++ implementation utilizing SIMD vector extensions and multi-threading. 

In a related vein, Nakandala \etal~\cite{nakandala2020hummingbird} proposed Hummingbird, which compiles decision trees into tensor operations to exploit GPU execution pipelines, demonstrating high speedups on large batch sizes. Conversely, Koschel \etal~\cite{koschel2023arm} examined ARM architectures, adapting the Quick\-Scorer bitvector traversal algorithm~\cite{lucchese2015quickscorer} to ARM NEON intrinsics and achieving up to $9.4\times$ speedups on IoT-class devices. Our work applies the fundamental principle of lockstep traversal directly within pure NumPy, eliminating C-extension compilation dependencies while matching cache hierarchies to maintain high throughput across both small and large batch regimes.

\subsection{Ensemble Early Exit and Cascaded Architectures}
Two-stage cascaded architectures---employing an ultra-lightweight preliminary filter to drop obviously benign traffic before invoking a heavier multi-class ensemble---are well established in embedded intrusion detection~\cite{ferreira2026canbus}. In parallel, early termination of tree ensembles has been explored to reduce computational cost: Hern\'andez-Lobato \etal~\cite{hernandezlobato2009pruning} prune trees dynamically based on statistical confidence intervals, while Wang \etal~\cite{wang2018qwyc} jointly optimize evaluation order and decision thresholds. However, these probabilistic early-stopping methods trade a small, non-zero error rate for execution speed. In an edge architecture governed by dual runtimes, any non-zero approximation inevitably causes cross-engine verdict divergence. Our proposed early-exit mechanism strictly enforces mathematical margin bounds, guaranteeing exact, zero-flip equivalence with the full ensemble by construction.

\subsection{Differential Verification and Model Parity Testing}
Differential testing across disparate execution engines is an active area of software engineering research. In deep learning, frameworks such as CRADLE~\cite{pham2019cradle} execute models across multiple backends to isolate internal compiler bugs, while Wen \etal~\cite{wen2025mindthegap} and Chen~\cite{chen2026integeralibi} investigate numerical divergences and bitwise inaccuracies across INT8 and quantized accelerator kernels. Cross-framework fuzzing~\cite{duan2026xamt} further uncovers divergent edge-case comparison semantics. 

For tree ensembles, symbolic verification tools such as Veritas~\cite{devos2021veritas} and sensitivity analysis methods~\cite{naik2026sensitivity} verify adversarial robustness over continuous real-valued bounds. However, they typically abstract away floating-point rounding realities. Official converter documentation (\eg \texttt{sklearn-\allowbreak onnx}) notes that converting double-precision splits to single-precision float32 can trigger label flips due to tree step-\allowbreak discontinuities~\cite{sklearnonnx_float}. Our research formalizes this phenomenon into an operational release gate, providing the first deterministic boundary probe generation algorithm that systematically audits runtime equivalence without relying on production traffic replay.
